\documentclass{article}
\usepackage[utf8]{inputenc}
\usepackage{hyperref}
\usepackage{geometry}
\geometry{a4paper, margin=1in}
\usepackage{xcolor}
\usepackage{listings}

\title{MLOps Project: Serving Role Implementation Report}
\author{}
\date{\today}

\begin{document}
\maketitle

\section{Introduction and Requirements}
The overarching requirement for the \textit{Serving Role} of this MLOps project was to prepare a comprehensive set of measured serving options that balance the trade-offs between fast, good, and cheap inference. As mandated by the project guidelines, this evaluation requires comparing multiple deployment approaches during system integration. 

Crucially, the course documentation emphasizes the necessity of measuring these metrics (latency, throughput, error rate) on realistic compute platforms. Thus, all evaluations are strictly intended to be executed on containerized \textit{Chameleon} instances rather than constrained local environments, explicitly adhering to the ``Right-sizing'' deployment evaluation paradigms. For the purposes of this milestone, we engineered the serving systems around the un-trained base architecture (the PANN-Cnn14 Audio Embedding Model).

\section{Implemented Serving Options: What and Why}

Drawing directly from Chapter 7 of the course notes (\textit{Model Serving}), we implemented a multi-stage optimization strategy evaluating the baseline, model-level, system-level, and infrastructure-level optimizations.

\subsection{Baseline Serving (FastAPI)}
\textbf{What was done:} We implemented a basic HTTP prediction endpoint utilizing the standard FastAPI web framework (\texttt{src/serving/baseline\_api.py}) packaged through \texttt{Dockerfile.baseline}.
\newline
\textbf{Why it was done:} As outlined in Section 7.4.1 (Designing the prediction service), deploying a canonical HTTP REST API is the simplest reference architecture. It provides an unoptimized control measure (yielding base P50/P95 latency and fundamental throughput limits) against which to benchmark further architectural optimizations.

\subsection{Model-Level Optimization (ONNX Runtime)}
\textbf{What was done:} We wrote \texttt{src/serving/export\_onnx.py} to freeze the PyTorch computational graph and compile it into an ONNX representation (\texttt{models/Cnn14.onnx}). We then consumed this optimized graph inside a dedicated FastAPI variant (\texttt{src/serving/onnx\_api.py}) executing via the \texttt{onnxruntime} \texttt{CPUExecutionProvider}.
\newline
\textbf{Why it was done:} Following Section 7.3.2 (Compiling and optimizing the graph), default PyTorch execution utilizes dynamic computational graphs which inherit substantial Python interpreter overhead. Compiling the model to an intermediate representation like ONNX allows for aggressive model-level operator fusion and memory pre-computation, consistently resulting in strictly lower latency for static batch sizes.

\subsection{System-Level Optimization (Triton Inference Server)}
\textbf{What was done:} We architected a dedicated \texttt{model\_repository} hierarchy focusing on a strict \texttt{config.pbtxt} definition. Specifically, we configured \texttt{dynamic\_batching} with a \texttt{max\_queue\_delay\_microseconds} of 1000 and preferred batch distributions. To simulate an end-to-end user request matching the other distinct APIs, we implemented an HTTP proxy via \texttt{src/serving/triton\_client.py}.
\newline
\textbf{Why it was done:} Section 7.4.3 emphasizes \textit{Dynamic Batching} to massively maximize hardware utilization (throughput) without sacrificing individual request latency. Triton uniquely intercepts incoming asynchronous network requests, queues them up to an optimal \texttt{preferred\_batch\_size}, and executes them utilizing vectorized tensor operations natively in C++. This acts as a definitive system-level optimization targeted specifically at heavy concurrent workloads that the baseline would drop.

\subsection{Bonus Infrastructure Framework (Ray Serve)}
\textbf{What was done:} We integrated the \texttt{ray.serve} deployment library (\texttt{src/serving/ray\_serve\_api.py}), dynamically scaling the model actor across configurable replica targets. We utilized the \texttt{@serve.batch} decorator to enable application-level queuing and batching, entirely enclosed within native Python. 
\newline
\textbf{Why it was done:} The project requirements explicitly outline a bonus for incorporating a serving framework not strictly covered in the lab (such as Ray Serve or KServe) that meaningfully improves design. Ray Serve improves upon standard Triton deployments by seamlessly permitting flexible, Pythonic multi-node autoscaling and complex orchestration without needing to maintain rigid C++ backed repository configurations. Section 7.4.4 directly addresses scaling with replicas---Ray Serve trivially achieves this robust horizontal scaling via its actor abstraction.

\section{Deliverables and Usage Guide}

\textbf{1. Model Artifacts and Serving Code:} The repository now encompasses distinct endpoints under \texttt{src/serving/}, ONNX export macros, and separated Dockerfiles corresponding to each optimization stage (\texttt{Dockerfile.baseline}, \texttt{Dockerfile.onnx}, \texttt{Dockerfile.triton\_client}, \texttt{Dockerfile.ray}).

\textbf{2. Evaluation Suite:} We built a heavily threaded load-testing utility (\texttt{scripts/benchmark\_serving.py}). It mocks authentic traffic by blasting concurrent asynchronous \texttt{POST} requests utilizing dummy WAV geometries against target endpoints, calculating deterministic throughput and P50/P95 distributions.

\textbf{3. Result Formatting:} Additionally provided is a Jupyter notebook (\texttt{notebooks/serving\_options\_analysis.ipynb}) designed to ingest the raw logging constants retrieved from the server, and render the final cleanly formatted Markdown \textit{Serving options table}.

\subsection{How to Execute (Chameleon Target)}
To produce the final table and sped-up demo video:
\begin{enumerate}
    \item Provision an active \texttt{compute\_icelake} bare-metal instance via the Chameleon GUI.
    \item SSH into the allocated node and \texttt{git clone} this project repository.
    \item Build and hydrate the desired container context: \texttt{docker build -t audio-serving-[type] -f Dockerfile.[type] .} 
    \item Spin up the instance: \texttt{docker run -p 8000:8000 audio-serving-[type]}.
    \item Open a concurrent terminal window and blast the instantiated endpoint representing usage: \texttt{python scripts/benchmark\_serving.py --url http://127.0.0.1:8000/predict -n 100 -c 10}
    \item Once complete, pipe those measured parameters back into the prepared Jupyter notebook configuration to export the ultimate grading table Markdown sheet, alongside attaching the screen recording of the benchmark execution phases.
\end{enumerate}

\end{document}
